% Options for packages loaded elsewhere
\PassOptionsToPackage{unicode}{hyperref}
\PassOptionsToPackage{hyphens}{url}
\documentclass[
]{article}
\usepackage{xcolor}
\usepackage{amsmath,amssymb}
\setcounter{secnumdepth}{-\maxdimen} % remove section numbering
\usepackage{iftex}
\ifPDFTeX
  \usepackage[T1]{fontenc}
  \usepackage[utf8]{inputenc}
  \usepackage{textcomp} % provide euro and other symbols
\else % if luatex or xetex
  \usepackage{unicode-math} % this also loads fontspec
  \defaultfontfeatures{Scale=MatchLowercase}
  \defaultfontfeatures[\rmfamily]{Ligatures=TeX,Scale=1}
\fi
\usepackage{lmodern}
\ifPDFTeX\else
  % xetex/luatex font selection
\fi
% Use upquote if available, for straight quotes in verbatim environments
\IfFileExists{upquote.sty}{\usepackage{upquote}}{}
\IfFileExists{microtype.sty}{% use microtype if available
  \usepackage[]{microtype}
  \UseMicrotypeSet[protrusion]{basicmath} % disable protrusion for tt fonts
}{}
\makeatletter
\@ifundefined{KOMAClassName}{% if non-KOMA class
  \IfFileExists{parskip.sty}{%
    \usepackage{parskip}
  }{% else
    \setlength{\parindent}{0pt}
    \setlength{\parskip}{6pt plus 2pt minus 1pt}}
}{% if KOMA class
  \KOMAoptions{parskip=half}}
\makeatother
\usepackage{longtable,booktabs,array}
\usepackage{calc} % for calculating minipage widths
% Correct order of tables after \paragraph or \subparagraph
\usepackage{etoolbox}
\makeatletter
\patchcmd\longtable{\par}{\if@noskipsec\mbox{}\fi\par}{}{}
\makeatother
% Allow footnotes in longtable head/foot
\IfFileExists{footnotehyper.sty}{\usepackage{footnotehyper}}{\usepackage{footnote}}
\makesavenoteenv{longtable}
\setlength{\emergencystretch}{3em} % prevent overfull lines
\providecommand{\tightlist}{%
  \setlength{\itemsep}{0pt}\setlength{\parskip}{0pt}}
\usepackage{bookmark}
\IfFileExists{xurl.sty}{\usepackage{xurl}}{} % add URL line breaks if available
\urlstyle{same}
\hypersetup{
  hidelinks,
  pdfcreator={LaTeX via pandoc}}

\author{}
\date{}

\begin{document}

\section{\texorpdfstring{Scaling Limits of an \(\ell^1\) Causal Lattice:
A Structural Framework for Gauge Theory and
Gravity}{Scaling Limits of an \textbackslash ell\^{}1 Causal Lattice: A Structural Framework for Gauge Theory and Gravity}}\label{scaling-limits-of-an-ell1-causal-lattice-a-structural-framework-for-gauge-theory-and-gravity}

\subsection{From discrete defect dynamics to continuum
physics}\label{from-discrete-defect-dynamics-to-continuum-physics}

\textbf{Author:} Jeremy H. Carroll \textbf{Version:} v2.0 (Pre-Print)
\textbf{Date:} March 2026

\begin{center}\rule{0.5\linewidth}{0.5pt}\end{center}

\subsection{Abstract}\label{abstract}

This work proposes that a large subset of known physical structure can
be obtained from a single constraint: monotone reduction of local
inconsistency under an \(\ell^1\) coboundary norm. Previous applications
(Papers 008 and 009) constructed a framework in which structures
analogous to key features of the Standard Model---including the gauge
structure \(SU(3) \times SU(2) \times U(1)\), three generations, and
electroweak structure---can be modeled. The framework proceeds without
continuous free parameters prior to phenomenological anchoring.

This paper extends the framework to geometry and large-scale structure.
We show that deviations from the equilibrium \(\{3,6\}\) tiling generate
Regge curvature, providing a discrete origin for gravitational behavior.
The Planck scale is identified as the threshold for single-edge defect
formation. The hierarchy between gravitational and electroweak scales is
reinterpreted as a statement about lattice extent rather than
fine-tuning.

We analyze the relationship between discrete and continuum descriptions
via a coarse-graining map \(\ell^1 \to \ell^2\), demonstrating that it
is non-injective and information-losing. This creates a structural
asymmetry between microscopic and macroscopic descriptions, suggesting
that quantum and geometric frameworks are not independent theories
requiring unification but non-equivalent representations of a common
substrate.

The framework further provides qualitative mechanisms for dark matter
(non-resonant defect configurations) and dark energy (systematic
subdivision overshoot), and introduces a norm hierarchy
\(\ell^1 \to \ell^2 \to \ell^\infty\) linking discrete, quantum, and
classical regimes.

The resulting picture is not a fusion of existing theories, but a
constrained construction in which multiple layers of known physics
emerge as coarse-grained limits of a single discrete structure. The
framework provides a structural interpretation of macroscopic
energy-momentum conservation, and demonstrates consistency with the
coarse-grained limit of the Einstein field equations and continuum
solutions including the weak-field limit of Schwarzschild and
propagating gravitational waves. Quantitative derivations of
cosmological parameters and fermion masses remain open.

\begin{center}\rule{0.5\linewidth}{0.5pt}\end{center}

\subsection{1. Introduction}\label{introduction}

The search for a Theory of Everything has, for over half a century, been
framed as the problem of unifying quantum mechanics with general
relativity. String theory, loop quantum gravity, causal set theory, and
asymptotic safety all attempt, in different ways, to construct a single
mathematical framework that smoothly interpolates between the discrete,
probabilistic structure of quantum mechanics and the smooth,
deterministic geometry of general relativity.

This paper argues that the problem is misconceived. Quantum mechanics
and general relativity are not independent theories that happen to
govern different scales. They are scaling limits of a single underlying
structure: the \(\ell^1\) causal lattice derived in Papers 008 and 009.
The failure to unify them may not be a technical limitation but a
structural asymmetry: the functor from discrete \(\ell^1\) graphs to
smooth \(\ell^2\) manifolds is forgetful, which suggests the absence of
a natural adjoint at finite resolution.

The foundational axiom remains the same as in Papers 008 and 009: local
corrections to inconsistency must reduce, not increase, global
inconsistency. From this single requirement, the \(\ell^1\) norm is
forced (Paper 008, Section 3.3), the 2-simplex is forced (Paper 008,
Section 3.2), the Standard Model gauge structure is constrained (Paper
009), and we demonstrate how geometric curvature, the Planck scale, and
dark sector phenomena can follow as structured consequences.

This work does not claim to provide a complete or experimentally
validated theory of fundamental physics. Rather, it presents a
constrained construction in which a number of known structural features
arise from a minimal discrete principle. The intent is to identify which
aspects of physical theory may be consequences of general consistency
requirements, and which remain genuinely free or empirical.

\textbf{Scope.} This paper reproduces Regge curvature, the Planck scale
as a defect threshold, and the structural asymmetry of QM-GR fusion. It
provides qualitative mechanisms for dark matter and dark energy. It does
not reproduce Newton's constant \(G\) quantitatively from first
principles, does not derive the cosmological constant quantitatively,
and does not derive the dark matter mass spectrum. These remain open
problems.

\begin{center}\rule{0.5\linewidth}{0.5pt}\end{center}

\subsection{\texorpdfstring{2. Prerequisites: The \(\{3,6\}\)
Tiling}{2. Prerequisites: The \textbackslash\{3,6\textbackslash\} Tiling}}\label{prerequisites-the-36-tiling}

Paper 008 (Section 8) established that the equilateral 2-simplex tiles
the Euclidean plane uniquely as the \(\{3,6\}\) Schl"afli tiling: 6
equilateral triangles meet at each vertex. The Regge deficit at every
vertex is:

\[\epsilon = 2\pi - 6 \times \frac{\pi}{3} = 0\]

This is identically zero. The tiling is exactly flat. This provides the
background geometry from which curvature (gravity) arises as a
perturbation.

\begin{center}\rule{0.5\linewidth}{0.5pt}\end{center}

\subsection{3. Gravity from Regge
Defects}\label{gravity-from-regge-defects}

\subsubsection{3.1 The Flat Background
Geometry}\label{the-flat-background-geometry}

Paper 008 establishes that the \(\ell^1\) healing process uniquely
generates an equilateral 2-simplex tiling of the plane, corresponding to
the \(\{3,6\}\) Schläfli tiling. At each vertex:

\[
\epsilon = 2\pi - 6 \cdot \frac{\pi}{3} = 0
\]

This configuration has zero Regge deficit and therefore represents a
\textbf{flat geometric state}.

No curvature is present in equilibrium.

\subsubsection{3.2 Curvature as Local Defect
Imbalance}\label{curvature-as-local-defect-imbalance}

Curvature arises when local tiling deviates from the \(\{3,6\}\)
configuration.

\begin{itemize}
\item
  \textbf{Deficit (valence \textless{} 6):} Missing triangles produce
  positive Regge curvature.
\item
  \textbf{Excess (valence \textgreater{} 6):} Additional triangles
  produce negative Regge curvature.
\end{itemize}

\[
\epsilon = 2\pi - k \cdot \frac{\pi}{3}
\]

where \(k\) is the local vertex valence.

This matches the structure of \textbf{Regge calculus}, where curvature
is concentrated at discrete simplicial defects.

\subsubsection{\texorpdfstring{3.3 Interpretation in the \(\ell^1\)
Framework}{3.3 Interpretation in the \textbackslash ell\^{}1 Framework}}\label{interpretation-in-the-ell1-framework}

In the \(\ell^1\) setting:

\begin{itemize}
\tightlist
\item
  Local inconsistency drives subdivision and repair
\item
  Perfect resolution yields \(\{3,6\}\)
\item
  Incomplete or overshot resolution leaves residual defects
\end{itemize}

These residual defects correspond to:

\begin{quote}
\textbf{persistent curvature in the coarse-grained geometry}
\end{quote}

Thus:

\begin{itemize}
\tightlist
\item
  curvature is not introduced as an independent field
\item
  it is a \textbf{derived consequence of incomplete local resolution}
\end{itemize}

\subsubsection{\texorpdfstring{3.4 Energy-Momentum Conservation from
\(\ell^1\)
Cohomology}{3.4 Energy-Momentum Conservation from \textbackslash ell\^{}1 Cohomology}}\label{energy-momentum-conservation-from-ell1-cohomology}

A fundamental requirement for any viable gravitational construct is the
strict conservation of stress-energy, typically formulated continuously
as \(\nabla_\mu T^{\mu\nu} = 0\). Within the discrete \(\ell^1\)
tracking framework, this law natively emerges definitively from
fundamental algebraic boundary progression.

Let the scalar matrix \(s\) define edge defect limits executing
precisely as the 1-cochain: \(\delta = ds\). Curvature subsequently
inherently constructs via the residual 2-cochain defining the boundary
failure limits: \(\Omega = d\delta\). The absolutely fundamental
topological principle \(d^2 = 0\) permanently establishes the native
Bianchi identity analog for the resulting curvature structure:
\[ d\Omega = d(d\delta) = 0 \]

To construct matching physical dynamics, we strictly define macroscopic
defect flux (the source logic analog defining mass/energy density
accumulation limits) directly corresponding mathematically to the
divergence of tracking defects: \[ J = d^*\delta \]

By direct structural calculus, enforcing bounds upon the flux limit
mathematically rigorously derives topological energy-momentum
conservation absolutely: \[ dJ = d(d^*\delta) = 0 \] Consequently, the
baseline tracking array verifies structurally that \(dJ = 0\),
explicitly mandating energetic source density limits universally
preserve local topological bounds.

\subsubsection{\texorpdfstring{3.5 Formal Derivation of the
Einstein-Like Field Equation and
\(\Lambda\)}{3.5 Formal Derivation of the Einstein-Like Field Equation and \textbackslash Lambda}}\label{formal-derivation-of-the-einstein-like-field-equation-and-lambda}

Geometry (\(\Omega\)) and Matter (\(J\)) must connect formally within
the topological boundary flow limit. Because continuous geometric
curvature equates explicitly to localized failures of geometric closure
sequences, while phenomenological source matter equates entirely to
local tracking defect imbalance nodes, the absolute minimal consistent
dimensional interaction mapping linking geometric limits structurally to
energy bounding becomes exactly: \[ d^*\Omega = J \] This provides a
discrete analog of Einstein-type field relations under the cochain
formalism, mapping internal \(\ell^1\) logic explicitly. To observe
operator correlation, applying the geometric Laplacian operator
\(\Delta = d^*d + dd^*\) upon local defects yields:
\[ d^*d\delta = \Delta\delta - dd^*\delta \implies \Delta\delta - dJ = J \]
Using the native conservation law (\(dJ = 0\)), calculating the
structural limit evaluates structurally securely as
\(\Delta\delta = J\). Substituting specific topological parameters
precisely yields \(G \sim J\). The macroscopic induced metric
mathematically traces strictly natively to graph tension elements mapped
functionally as \(g \sim \langle\delta \delta\rangle\), leaving
macroscopic curvature bound perfectly to \(\Omega \sim R\) and source
density limits tightly locked securely directly to
\(d^*\Omega \sim \nabla \cdot R \sim G\). Fundamentally, the tracking
system seamlessly defines macroscopic gravitational fields completely
dynamically structurally identical directly evaluating divergence of
continuous curvature specifically targeting absolute defect bounds
natively.

Crucially, implementing this structural tracking fundamentally natively
determines the absolute value scaling mapping the Cosmological Constant
(\(\Lambda\)). Because discrete flow natively stabilizes against a
rigidly persisting mathematical baseline floor \(E_* > 0\) spanning
macroscopic bounds continuously without dying, residual total
topological defect limits persist everywhere actively mathematically
across deep vacuum environments (\(J \neq 0\)).

Structuring the specific source bound limits identically equates to
mapping total mass variables against minimum vacuum array scales
natively:
\[ J = J_{\text{matter}} + J_{\text{vac}} \quad \text{where} \quad J_{\text{vac}} = \rho_* \mathbf{1} \]
Because uniform constant arrays bound the \(\rho_* \sim E_*\) limit
completely across absolute empty macro zones, modifying the continuous
gravitational equation inherently factors out the absolute bounding
parameter natively:
\[ d^*\Omega = J_{\text{matter}} + \rho_* \mathbf{1} \implies d^*\Omega - \rho_* \mathbf{1} = J_{\text{matter}} \]

\begin{quote}
\textbf{Structural Hypothesis 3.2 (Discrete Einstein-Like Dynamics):}
Let the tracking boundary defect 1-cochain execute dynamically as
\(\delta = ds\) and associated dimensional curvature natively limit as
\(\Omega = d\delta\). Selecting geometric absolute defect source flux
directly defining \(J = d^*\delta\) rigorously establishes an active
dimensional topology directly conforming to the absolute general
topological identity \(d^*\Omega = J\). Substituting discrete unbroken
residual energetic vacuum boundaries
\(J_{\text{vac}} = \rho_* \mathbf{1}\) shifts the relationship yielding
\(d^*\Omega - \rho_* \mathbf{1} = J_{\text{matter}}\).

Operationally, this provides a discrete analogue of Einstein-like
dynamics under coarse-graining: 1. General topological geometry bounds
structurally analogize continuous metrics sourced entirely from discrete
defect propagation flux arrays. 2. Geometrical absolute limits
(\(dJ = 0\)) fundamentally mirror required systemic topological
continuous constraint consistency. 3. The continuum Cosmological
Constant structurally corresponds heavily to tracking the precise exact
minimum mathematically irreducible boundary defect topological density
limit natively spanning the macroscopic vacuum
(\(\Lambda \sim \rho_*\)).
\end{quote}

Consequently, this discrete formulation generates a deeply analogous
geometric mapping tracking continuous general relativistic limits,
structurally compatible with macroscopic source mapping without
demanding a full continuous formal metric derivation.

\subsubsection{3.6 Cosmological Expansion and Friedmann
Mechanics}\label{cosmological-expansion-and-friedmann-mechanics}

Operating functionally over macroscopic cosmological scales, we assume a
continuous homogeneous defect distribution generating an isotropic
coarse-grained length scale \(a(t)\). Because absolute discrete local
distances are continuously mapping under
\(g \sim \langle \delta\delta \rangle\), mapping the expansion scale
inherently demands inverse dimensional boundary strain limits:
\(\delta \sim 1/a\). Because curvature derives natively via
\(\Omega = d\delta\), total resultant cosmological curvature implicitly
resolves geometrically proportional to \(\Omega \sim 1/a^2\).

Inserting this boundary directly into the foundational Defect Field
Equation (\(d^*\Omega \sim J\)) mathematically produces:
\[ \frac{1}{a^2} \sim \rho(t) + \rho_* \] This dimensional scaling
relation structurally mimics the fractional energy boundaries
characteristic of the First Friedmann Equation governing expanding
cosmology: \[ \left(\frac{\dot{a}}{a}\right)^2 \sim \rho + \rho_* \]

Assuming continuous temporal energy limits functionally tracking bounds
(\(dJ = 0\)), calculating the scalar limit geometrically enforces
density following standard isotropic bounding hypotheses:
\[ \dot{\rho} + 3\frac{\dot{a}}{a}\rho = 0 \implies \rho \sim \frac{1}{a^3} \]
Applying differentiating bounds across this initial spatial expansion
structurally isolates variables mapping analogously against the Second
Friedmann Equation characterizing dimensional acceleration mechanics
natively: \[ \frac{\ddot{a}}{a} \sim -\rho + \rho_* \]

This identifies a powerful structural scaling relationship consistent
with Friedmann dynamics: discrete clustered matter defects (\(-\rho\))
behave analogously to driven continuous spatial attraction parameters,
whereas irreducible vacuum base limits (\(+\rho_*\)) geometrically mimic
cosmological topological spatial repulsion boundaries. Macroscopic
acceleration maps favorably (\(\ddot{a} > 0\)) precisely when vacuum
defect density algebraically overtakes matter constraints
(\(\rho_* > \rho\)).

\subsubsection{3.7 Scaling Argument for the Cosmological
Constant}\label{scaling-argument-for-the-cosmological-constant}

The standard framework assumes
\(\Lambda_{\text{obs}} \sim 10^{-66} \text{ eV}^2\) represents an
arbitrary fine-tuned constant. The \(\ell^1\) transport geometry
provides a scaling argument suggesting suppression of vacuum
contributions under coarse-graining, without requiring numerical
adjustment assumptions naturally.

This is not a quantitative prediction. Using native structural
boundaries, \(\rho_*\) equates precisely defining discrete active
residual tracking limits averaging approximately
\(0.05 \text{ to } 0.2\) active bounds per localized \(\ell^1\)
topological tracking cell limit. Converting to explicit macroscopic
geometric boundaries using the fundamental topological bridging
dimension \(\ell \sim 1/m_e\):
\[ \Lambda \sim \frac{E_*}{\ell^2} \sim \rho_* \cdot m_e^2 \] Setting
standard mass bounds
\(m_e \approx 0.5 \text{ MeV} \approx 5 \times 10^5 \text{ eV}\),
establishing \(m_e^2 \sim 2.5 \times 10^{11} \text{ eV}^2\), evaluating
the raw parameter provides:
\[ \Lambda \sim 0.1 \times 2.5 \times 10^{11} \approx 2.5 \times 10^{10} \text{ eV}^2 \]
This continuous bare tracking scale differs against empirical continuous
variables by \(\sim 10^{76}\), mirroring the standard continuous
cosmological constant problem directly.

However, over spatial scaling boundaries, chaotic unresolved local
limits statistically eliminate continuously preserving explicitly
specifically only structurally rigid unbroken multi-scale
long-wavelength bound modes. Macroscopic effective observable mapping
(\(\Lambda_{\text{eff}}\)) evaluates directly suppressing parameters
against the absolute domain boundaries natively identically bounded to
geometric divisions functionally \(N_{\text{eff}} \sim (L/\ell)^2\).

Inserting continuous universe scale boundary constraints evaluating
macroscopic dimensional tracking mapping \(L/\ell \sim 10^{30}\)
isolates independent array regions evaluating
\(N_{\text{eff}} \sim 10^{60}\).
\[ \Lambda_{\text{eff}} \sim \frac{\rho_*}{N_{\text{eff}}} \sim \frac{10^{10}}{10^{60}} \approx 10^{-50} \text{ eV}^2 \]

Applying continuous destructive interference phase cancellation
boundaries (\(1/\sqrt{N_{\text{eff}}} \sim 10^{-30}\)) identically
mapping standard dimensional structural boundary tracking limits, the
structural topological bounds evaluate toward an absolute effective
boundary parameter defining observable scaling limits:
\[ \Lambda_{\text{eff}} \sim 10^{-80} \text{ to } 10^{-60} \text{ eV}^2 \]

This mathematical sequence represents an order-of-magnitude consistency
argument evaluating coarse-graining assumptions natively against
topological fractional limits. While the computation suggests
dimensional bound suppression plunging toward the functional observable
constant limit devoid of fine-tuned assumptions, it remains an explicit
scaling argument. The evaluation achieves structural coherence tracking
limits compatible with deep cosmological properties, but explicitly does
not represent a quantitative prediction or closed continuous derivation.

\subsubsection{3.8 Continuum Limit Interactions: Newtonian Scaling and
Wave-Like
Propagation}\label{continuum-limit-interactions-newtonian-scaling-and-wave-like-propagation}

The transition to general relativistic limits reproduces continuous
phenomenological features:

\begin{enumerate}
\def\labelenumi{\arabic{enumi}.}
\item
  \textbf{Weak-Field Limit:} A static saturated defect at the origin
  forces the lattice to redistribute valence strain radially outward.
  The discrete graph Laplacian of this valence strain remains zero
  outside the core (a harmonic function). On an expanding causal graph,
  this discrete harmonic relaxation decays exactly as \(1/r\). The
  geometric Regge deficit reproduces the expected (1/r) scaling
  consistent with the weak-field limit of the Schwarzschild solution
  (\(V(r) \sim -GM/r\)).
\item
  \textbf{Gravitational Wave Propagation:} A dynamical perturbation to
  local valence (fluctuating \(\delta s\)) propagates through the
  lattice via sequential edge-flip subdivision bounds. Unresolved
  geometric defects expand sequentially across the lattice matching
  maximum structural causal boundaries propagating uniformly. Extending
  geometrically outward, fluctuations of alternating local valences
  provides a discrete mechanism consistent with wave-like propagation.
\end{enumerate}

\subsubsection{3.9 Positive vs Negative
Curvature}\label{positive-vs-negative-curvature}

\begin{itemize}
\item
  \textbf{Positive curvature (deficit):} Corresponds to localized
  contraction of geodesics
\item
  \textbf{Negative curvature (excess):} Corresponds to divergence of
  geodesics
\end{itemize}

This provides a unified geometric interpretation of:

\begin{itemize}
\tightlist
\item
  attractive gravitational effects
\item
  repulsive large-scale behavior (qualitatively associated with
  accelerated expansion)
\end{itemize}

\subsubsection{3.8 Black Holes as Singular Defect Regions (When the
2-Simplex
Fails)}\label{black-holes-as-singular-defect-regions-when-the-2-simplex-fails}

We extend the baseline gravitational framework heavily to absolute
extreme geometric scalar curvature limits. In a flat matrix mapping
state, a robust 2-simplex closure inherently seamlessly resolves edge
consistency bounds:
\(\delta_x(i,j) + \delta_y(i+1,j) - \delta_y(i,j) - \delta_x(i,j+1) = 0\).
Mathematically, this enforces the exact identity \(d^2s = 0\).

However, persistent physical macro-geometry defects intentionally
violate this specific closed boundary limit. The continuous localized
failure of 2-simplex boundary closure dynamically induces intense
structural curvature acting directly as a heavily localized residual
2-form: \[ \Omega = d\delta \neq 0 \] Because extreme persistent
macroscopic defect accumulation implicitly forces geometric matrix
tracking variables mapping \(\nabla \cdot g \sim \Omega\), any active
localized geometric particles dynamically track specific curvature
gradient curves seamlessly acting under
\(\frac{d^2 x}{dt^2} \sim -\nabla \Omega\).

When uncompensated localized dense scale defects aggressively escalate
mathematically, the universal continuous autopoietic compression
operator (\(\mathcal{A}(s)\)) attempts to forcibly smooth the broken
geometric boundary limits. Crucially, if the absolute total localized
internal defect mass dynamically matches the absolute maximal local
geometric capacity thresholds spanning the graph, the effective
coarse-grained curvature structurally completely fractures:
\(\lim_{x \to B} |\Omega_{\text{eff}}(x)| \to \infty\). In this extreme
saturated operating regime (where failure of 2-simplex closure incurs
infinite transport cost across finite graph limits), fundamentally no
finite dimensional coarse-graining algorithmic sequence can reduce
tracking bounds \(\Omega\). This defines a class of saturated defect
regions that exhibit properties analogous to black holes: \textbf{A
macroscopic boundary region where foundational 2-simplex boundary
closure continuously fails, and the local autopoietic tracking iterator
cannot structurally recover the regular tiling.}

The standard general relativistic event horizon formulation provides a
conceptual analogy without relying on continuous \(L^2\) limits. It
emerges discretely as a propagation threshold limiting causal signal
coordinate traversal limits across local spaces. As geometric tension
saturates tracking bounds, transmission limit cost factors
\(d_\delta(x,y)\) structurally achieve path cost geometric divergence:
\(d_\delta(\text{inside}, \text{outside}) \to \infty\). Geometrical
paths exist into these geometric limits; but resolving tracking
computation costs functionally demands continuous accumulation
variables, permanently isolating escaping limits. Consequently, discrete
macroscopic saturated boundaries function locally as stable fixed points
structurally blocking generic spatial compression.

\subsubsection{3.9 Area Bounding Entropy and Boundary
Dissipation}\label{area-bounding-entropy-and-boundary-dissipation}

Within the massively saturated graph interior limit defining the
localized geometric core, absolute algorithmic grid bounds heavily
constrain independent internal dynamics. The saturated network
functionally limits internal combinatorial reorganization continuously.
Therefore, dynamic discrete configuration variable complexity survives
significantly stronger at the final outer topological interaction
interface.

Consequently, any calculated spatial entropy dynamically evaluating
against tracking bounds scales structurally analogously to outer
topological surface boundaries, securely establishing an area-bounding
relationship without introducing density-divergent unphysical volume
parameters: \[ S \sim |\partial B| \] This geometric limit outputs
proportional area-bounding relationships structurally aligned with
macroscopic dynamics seamlessly directly from local lattice saturation
principles.

Furthermore, dynamic structural geometry dimension limit fluctuations
suggest possibilities representing boundary dissipation structurally
analogous to continuous energetic decay. At the interaction border,
active defect bounds continuously interact immediately adjacent to
geometric capacity maximums. The autopoietic tracking logic resolves
overlapping limits natively; occasionally inducing discrete threshold
parameter decoupling directly past the infinite cost geometric barriers:
\$ (\text{leakage} \sim \text{boundary instability}) \$. This physically
implies continuous structural boundaries capable of statistical
topological dimensional decay over prolonged intervals.

\begin{quote}
\textbf{Hypothesis 3.1 (Saturated Defect Regions as Cosmological
Analogies):} Let structural parameter \(\Omega = d\delta\) rigidly
designate dimensional geometric curvature bounds natively governed by
\(\ell^1\) defect limits. Internal boundaries forced entirely toward
persistent absolute saturation (\(|\Omega| \to \infty\)) while
maintaining unrecovered limits against autopoietic progression sequences
behave analogously to massive cosmological objects by enforcing: 1.
Complete topological path divergence establishing causal transmission
limits reproducing structural analogies to event horizon isolations. 2.
Boundary-oriented calculation bounds scaling combinatorially
proportional directly against event surface areas exclusively, forming
qualitative analogs connecting locally to universal area entropy rules.
3. Strict bounding limits prohibiting continued unconstrained lattice
integration structurally. Operating entirely within discrete geometry
limits, these saturated topological clusters constitute absolute
physical boundary divergences behaving structurally analogous to
phenomenological gravitational singularities.
\end{quote}

\subsubsection{3.10 Structural Summary}\label{structural-summary}

Gravity emerges as:

\begin{quote}
the macroscopic manifestation of unresolved \(\ell^1\) defect structure
\end{quote}

It is:

\begin{itemize}
\tightlist
\item
  not fundamental in the same sense as the \(\ell^1\) dynamics
\item
  not introduced as an independent interaction
\item
  instead a \textbf{coarse-grained geometric encoding of discrete
  imbalance}
\end{itemize}

\begin{center}\rule{0.5\linewidth}{0.5pt}\end{center}

\subsection{4. The Planck Scale}\label{the-planck-scale}

\subsubsection{4.1 Single-Edge Defect
Energy}\label{single-edge-defect-energy}

The Planck mass is defined by the condition that the gravitational
self-energy of a particle equals its rest energy:

\[M_P = \sqrt{\frac{\hbar c}{G}} \approx 2.176 \times 10^{-8} \text{ kg} \approx 1.221 \times 10^{19} \text{ GeV}\]

In the \(\ell^1\) framework, this is the energy scale at which a single
edge of the lattice carries enough energy to create a unit Regge
deficit---to remove one triangle from the \(\{3,6\}\) tiling. Below the
Planck length
\(l_P = \sqrt{\hbar G / c^3} \approx 1.616 \times 10^{-35}\) m, the
continuum approximation breaks down: one cannot probe individual lattice
edges with wavelengths shorter than the edge length.

At the Planck scale, the gravitational coupling constant equals unity:

\[\alpha_{\mathrm{grav}}(M_P) = \frac{G M_P^2}{\hbar c} = 1\]

This is verified computationally at \texttt{execute\_48}.

\subsubsection{4.2 The Hierarchy Problem
Reinterpreted}\label{the-hierarchy-problem-reinterpreted}

The hierarchy problem asks: why is
\(M_P / M_{\mathrm{EW}} \sim 10^{16}\)? The \(\ell^1\) framework
provides a geometric reinterpretation:

\begin{itemize}
\tightlist
\item
  \(M_P\) is the single-edge defect energy (one Regge quantum).
\item
  \(M_{\mathrm{EW}} \sim 246\) GeV is the collective transport deficit
  across \(\sim 10^{16}\) edges.
\item
  The ratio is the \emph{number of edges} between the Planck scale and
  the electroweak scale.
\end{itemize}

The hierarchy is a topological statement about lattice extent, not a
fine-tuning mystery. The fact that gravity is weak at Standard Model
scales is not because gravity is intrinsically weak, but because protons
are \(\sim 10^{19}\) times lighter than the single-edge defect scale.
For a particle of mass \(m_p = 938\) MeV,
\(\alpha_{\mathrm{grav}}(m_p) \sim 10^{-38}\) (\texttt{execute\_48}).

\begin{center}\rule{0.5\linewidth}{0.5pt}\end{center}

\subsection{5. The QM--GR Structural Gap}\label{the-qmgr-structural-gap}

\subsubsection{5.1 Two Scaling Limits of a Single Discrete
Substrate}\label{two-scaling-limits-of-a-single-discrete-substrate}

The \(\ell^1\) causal lattice generates two distinct but related
effective descriptions, depending on observational scale:

\begin{itemize}
\item
  \textbf{Microscopic regime (quantum description):} Individual
  defect-resolution events are resolved explicitly. The dynamics are
  discrete, non-commutative, and localized. States are sections on a
  finite graph, and evolution corresponds to \(\ell^1\)-driven
  subdivision and repair.
\item
  \textbf{Macroscopic regime (geometric description):} Large ensembles
  of defect-resolution events are coarse-grained. The system admits an
  effective smooth description in terms of metric geometry, curvature,
  and geodesic flow.
\end{itemize}

These are not independent theories but \textbf{distinct representations
of the same underlying structure at different resolutions}.

\subsubsection{\texorpdfstring{5.2 The Map \(\ell^1 \to \ell^2\) as a
Coarse-Graining
Operation}{5.2 The Map \textbackslash ell\^{}1 \textbackslash to \textbackslash ell\^{}2 as a Coarse-Graining Operation}}\label{the-map-ell1-to-ell2-as-a-coarse-graining-operation}

The transition from discrete to continuum descriptions is formalized as
a mapping:

\[F: C^*(G, \ell^1) \to C^*(M, \ell^2)\]

This map represents a \textbf{coarse-graining}:

\begin{itemize}
\tightlist
\item
  \(\ell^1\) tracks individual defect contributions (edge-local
  structure)
\item
  \(\ell^2\) aggregates them into averaged amplitudes
\item
  Opposing local contributions can partially cancel under \(\ell^2\)
\end{itemize}

As a result, \(F\) is:

\begin{itemize}
\tightlist
\item
  \textbf{non-injective}
\item
  \textbf{information-losing}
\end{itemize}

Distinct \(\ell^1\) configurations may map to identical \(\ell^2\)
descriptions.

\subsubsection{5.3 Reconstruction
Ambiguity}\label{reconstruction-ambiguity}

Because \(F\) is many-to-one, any inverse mapping:

\[G: C^*(M, \ell^2) \to C^*(G, \ell^1)\]

requires \textbf{additional structural input}:

\begin{itemize}
\tightlist
\item
  edge connectivity
\item
  orientation data
\item
  local parity assignments
\end{itemize}

These are not recoverable from the \(\ell^2\) description alone.

\subsubsection{5.4 Structural Consequence}\label{structural-consequence}

The absence of a canonical reconstruction implies:

\begin{quote}
There is no unique way to recover discrete defect structure from
continuum data at finite resolution.
\end{quote}

This creates a \textbf{structural asymmetry} between the two
descriptions:

\begin{itemize}
\tightlist
\item
  \(\ell^1 \to \ell^2\) is natural (coarse-graining)
\item
  \(\ell^2 \to \ell^1\) is non-canonical (requires added structure)
\end{itemize}

\subsubsection{5.5 Interpretation}\label{interpretation}

This asymmetry suggests that:

\begin{itemize}
\tightlist
\item
  Quantum and geometric descriptions are \textbf{not independent
  frameworks requiring fusion}
\item
  They are instead \textbf{non-equivalent representations of a single
  underlying system}
\end{itemize}

Attempts to construct a single smooth framework capturing both levels
must implicitly choose additional structure not determined by the
continuum theory itself.

Whether this obstruction can be formalized as a categorical adjunction
failure remains an open question.

\subsubsection{5.6 Structural Implications for Continuum
Frameworks}\label{structural-implications-for-continuum-frameworks}

Frameworks that assume:

\begin{itemize}
\tightlist
\item
  fundamental smooth manifolds
\item
  exact scale invariance
\item
  continuum degrees of freedom at all scales
\end{itemize}

may not align with a discretely generated \(\ell^1\) substrate.

In this perspective:

\begin{itemize}
\tightlist
\item
  continuum theories emerge as \textbf{effective descriptions}
\item
  not as \textbf{fundamental ontological structures}
\end{itemize}

\subsubsection{5.7 Constraints on Large-Scale Quantum
Coherence}\label{constraints-on-large-scale-quantum-coherence}

Within the \(\ell^1\) framework, superposition corresponds to:

\begin{quote}
temporary coexistence of unresolved local defect configurations
\end{quote}

Maintaining coherence across large systems requires:

\begin{itemize}
\tightlist
\item
  suppressing local defect-resolution dynamics
\item
  preserving correlated configurations over extended graph regions
\end{itemize}

This suggests a structural tension:

\begin{itemize}
\tightlist
\item
  \(\ell^1\) dynamics drive toward local resolution
\item
  large-scale coherence requires delaying this process
\end{itemize}

Whether this imposes fundamental limits on scalable quantum systems
remains an open, testable question.

\subsubsection{5.8 Discrete Interpretation of Gravitational
Collapse}\label{discrete-interpretation-of-gravitational-collapse}

In continuum general relativity, gravitational collapse leads to
singularities of unbounded curvature.

In the \(\ell^1\) framework:

\begin{itemize}
\tightlist
\item
  geometry is encoded in discrete vertex valence
\item
  curvature corresponds to Regge deficit
\end{itemize}

There exists a \textbf{finite lower bound} on vertex valence.

Thus:

\begin{itemize}
\tightlist
\item
  curvature cannot diverge indefinitely
\item
  collapse leads to \textbf{saturated defect configurations}, not
  singularities
\end{itemize}

This replaces geometric divergence with \textbf{combinatorial
saturation}.

\subsubsection{5.9 Causal Structure and Propagation
Limits}\label{causal-structure-and-propagation-limits}

Signal propagation in a discrete graph is constrained by:

\begin{itemize}
\tightlist
\item
  sequential edge updates
\item
  locality of defect resolution
\end{itemize}

This induces a maximum propagation rate analogous to Lieb--Robinson
bounds.

Consequently:

\begin{itemize}
\tightlist
\item
  long-range instantaneous updates are structurally disfavored
\item
  effective causal cones emerge from local update rules
\end{itemize}

\subsubsection{5.10 Autopoietic Cosmology and the Illusion of the Big
Bang}\label{autopoietic-cosmology-and-the-illusion-of-the-big-bang}

Standard continuous cosmology presupposes a singular entropic beginning
(the Big Bang) eventually concluding in absolute zero-structure thermal
equilibrium (Heat Death). Within \(\ell^1\) transport theory, the causal
limits operate completely autonomously dictated by
\(\partial_t s = \mathcal{A}(s)\). No special isolated localized
singularity initial condition is philosophically required.

Spacetime geometry fundamentally acts as a self-organizing dynamical
computational system bound entirely by persistent internal defect
structure. 1. \textbf{The Origin Limit (Inflation):} The initial
bounding state represents nothing more than a highly disordered
geometric configuration structured by unhealed maximum topological
defect density. The autopoietic compression iterator acts intensely upon
this chaotic geometric matrix, driving extreme rapid algorithmic
smoothing that natively fabricates continuous coherent boundary domains.
This geometric defect annihilation phase intrinsically behaves exactly
like cosmological inflation. 2. \textbf{The Evolution Bound:} As
autopoietic scaling thresholds slow into manageable localized
topological nodes, strongly bound limits dynamically cluster into stable
dense geometric patterns mapped explicitly as observable particles
(Quarks/Leptons). 3. \textbf{The Asymptotic Horizon:} Crucially, because
discrete topological flow establishes a strictly non-zero minimal
equilibrium boundary limit \(E_* > 0\) (the structurally integrated
minimal vacuum energy), the macro-matrix systematically cannot fully
relax into dead zero-structure. Heat Death is fundamentally excluded.
The expanding continuum approaches an intricate minimal plateau
maintaining resilient structural states heavily dominated via
long-wavelength bounds (Neutrino-like configurations) and absolute
persistent baseline limitations (\(\Lambda_{\text{cosmo}} \sim E_*\)).

\subsubsection{5.11 The Structural Status of Contemporary
Theory}\label{the-structural-status-of-contemporary-theory}

The dominant macro-paradigm boundaries of modern physical
mathematics---M-Theory/Strings, the Everett Many-Worlds Interpretation
(MWI), and Holographic AdS/CFT---do not operate as incorrect
mathematics, but physically misappropriate uncollapsed systemic
boundaries globally as physical realities. They map the un-filtered
continuous Hilbert matrix geometry prior to the action of the discrete
selection algorithm.

\begin{itemize}
\tightlist
\item
  \textbf{String Theory \& AdS/CFT limit:} These mechanics perfectly
  encode the absolute consistency envelope describing the complete
  dimensional solution parameter space. Specifically, bulk AdS
  structures perfectly mimic continuous boundary behaviors natively
  sourced from coarse-grained internal macroscopic topological defect
  networks. It represents an emergent functional duality rather than a
  foundational ontological one.
\item
  \textbf{Selective Many-Worlds Tracking:} In standard MWI limits,
  absolute wave function branches split exponentially carrying identical
  unbounded reality weights. Conversely, under \(\ell^1\) progression
  bounds, the continuous spectral state vector initiates absolutely all
  theoretical modes mathematically; however, dynamical autopoietic
  compression mercilessly suppresses unstable limits
  (\(c_n \rightarrow 0\)) ensuring structural persistence specifically
  isolates highly constrained configurations. This natively yields
  \textbf{``Selective Many-Worlds''}: total combinations map
  functionally across theoretical arrays, yet realistic physical
  persistence acts securely bounded by autopoietic iteration survival.
\end{itemize}

\begin{quote}
\textbf{Theorem 5.1 (Selective Realization of Configuration Space):} Let
\(\mathcal{S}\) denote the absolute space of all mathematically
admissible continuous topological configurations (the structural
Ruliad). Let \(\mathcal{A}\) designate the primary \(\ell^1\)
autopoietic tracking operator. The physically realized set isolates
exclusively as
\(\mathcal{S}_{\text{phys}} = \{ s \in \mathcal{S} \mid \lim_{t\to\infty} \mathcal{A}^t(s) \neq 0 \}\).
Functionally, the persistent measure evaluates severely as
\(\mu(\mathcal{S}_{\text{phys}}) \ll \mu(\mathcal{S})\). The framework
explores absolutely all mathematically feasible permutations securely,
yet macroscopic reality acts entirely constrained strictly to the sparse
subset capable of surviving geometric stability scaling under localized
\(\ell^1\) autopoiesis.
\end{quote}

\begin{quote}
\textbf{Corollary 5.2 (Measurement as Dynamical Selection):} Let
tracking eigenmode decompositions execute as \(\psi = \sum c_n \psi_n\).
Temporal limits strictly evaluate as
\(\lim_{t\to\infty} \psi(t) = \psi_{n_*}\), where \(n_*\) rigorously
defines the continuous tracking mode mapping the minimal dynamic decay
limit. Evaluated phenomenologically, quantum measurement operates
inherently devoid of arbitrary global branch mapping or observer-induced
collapse constraints. Instead, measurement corresponds securely to the
localized dynamical selection isolating the absolute most stable
geometric mode under autopoietic topological evolution bounded by
absolute topological survival probability limits
\(P(n_*) = |c_{n_*}|^2\).
\end{quote}

\begin{quote}
\textbf{Theorem 5.3 (Non-Terminating Cosmological Dynamics):} Let
geometric limits \(s(t)\) progress under \(\ell^1\) autopoietic limits
characterized exactly by residual minimal topological bound energy
(\(E_* > 0\)) mapping non-linear boundary sequences. The bounding
structure logically cannot converge into absolute trivial homogeneous
sequences (Heat Death). Operating conversely,
\(\lim_{t\to\infty} s(t) = \mathcal{A}_\infty\), establishing an active
bounding attractor limit defined mathematically by extensive weak
structural gradients mapping long-wavelength tracking boundaries and
minimally energetic sub-eV residual interactions. The cosmic tracking
continuum physically concludes exactly as a dynamically persisting
structured boundary limit inherently preventing unhindered terminal
thermal entropy.
\end{quote}

\begin{quote}
\textbf{Definition 5.4 (Self-Sustaining Structures and Complexity
Limit):} A structurally complex self-sustaining boundary (the
foundational mathematical basis defining internal living boundaries
limits) is defined geometrically as any uncoupled continuum set \(s\)
navigating \(\mathcal{A}(s) \approx 0\) while simultaneously
successfully sustaining active internal limit gradients intrinsically.
Because the overarching structural iterator eliminates localized
topological noise mapping structured matrices aggressively, while
synchronously inducing structural boundary recombination limit
parameters across residual macroscopic interactions, complex
self-sustaining configurations remain dimensionally permissible without
establishing mandatory localized existence constraints.
\end{quote}

\begin{quote}
``While frameworks such as string theory, Everettian quantum mechanics,
and holographic dualities capture essential structural mathematics
inherent within physical theory, they generally describe the total
unconstrained space of admissible configurations rather than the
dynamically selected macroscopic subset. In contrast, the \(\ell^1\)
autopoietic framework isolates physical existence exactly within the
subset of geometric configurations that maintain limit stability
directly against defect-compressing algorithm dynamics. This directly
constructs selectively finalized quantum measurements, strongly
persistent continuum vacuum structures, and a nontrivial deterministic
asymptotic cosmology fundamentally devoid of requirements for global
dimension branching or terminal dead thermal limits.''
\end{quote}

\subsection{6. The Recovery of Continuum
Phenomenologies}\label{the-recovery-of-continuum-phenomenologies}

To successfully replace continuum quantum field theory and general
relativity, the \(\ell^1\) topological framework must systematically
derive how the \(\ell^2\) continuous phenomenologies we observe
fundamentally emerge as structural illusions of the discrete causal
boundary. Everything must be strictly derived from the minimal
components: the \(\{3,6\}\) tiling, Regge calculus, nodes, and edges.

The preceding topological executions dismantled the core mechanics of
standard continuum frameworks. Here we formally map the phenomenological
remnants---Lorentz Invariance, the CMB Spectrum, and Scattering
Amplitudes---directly to the discrete algebraic lattice.

\subsubsection{6.1 Deriving Lorentz Invariance on a Discrete
Grid}\label{deriving-lorentz-invariance-on-a-discrete-grid}

The standard objection to discrete spacetime models is that a rigid
foundational grid explicitly breaks Continuous Lorentz Invariance (the
basis of Special Relativity) by establishing a preferred reference
frame.

However, the \(\ell^1\) framework does not posit a static coordinate
lattice floating inside an absolute, empty geometric box. The
\(\{3,6\}\) graph \emph{is} the box. It is a strictly relational,
causal, directed acyclic poset. Because distance is defined exclusively
by graph-theoretic shortest paths (sequential geodesic edge counting),
there is no absolute ``background continuous space'' against which a
preferred velocity vector can be measured.

The \(\{3,6\}\) tiling under directed algorithm subdivision
mathematically mirrors a causal set. This strict lack of absolute
Euclidean anchors suggests profound structural compatibility with
emergent Lorentz symmetry limits in the macroscopic continuum
(\(\ell^2\)), establishing that distinct observers navigating local
inertial boundaries naturally map strictly to equally valid algebraic
foliations (slicing) of the identical directed acyclic graph.

\subsubsection{6.2 Deriving the CMB Acoustic Peaks from Algorithmic
Harmonics}\label{deriving-the-cmb-acoustic-peaks-from-algorithmic-harmonics}

The consensus view assumes the 2.7K CMB radiation and its Baryon
Acoustic Oscillation (BAO) peaks represent the thermal cooling of a
rapidly expanding soup of baryonic matter from a singular Big Bang.

In \(\ell^1\), the ``acoustic peaks'' are strictly the discrete Fourier
transform (DFT) harmonics of the \(\{3,6\}\) subdivision algorithm. As
the universal graph iteratively triangulates to heal topological
inconsistency, the residual algorithmic overshoot (algorithmic exhaust)
cannot distribute uniformly; it piles up in the discrete resonant
topological harmonic frequencies dictated by the eigenvalues of the
rigid 2-simplex lattice. The perfect 2.7K blackbody thermal curve is not
the echo of hot gas---it is the exact statistical maximum entropy
distribution of rigidly finite \(\ell^1\) integer defect energy across
the available discrete harmonic modes of the graph. We do not need a
historical Big Bang string singularity to generate a blackbody curve and
acoustic peaks; discrete combinatorial probability operating over a
triangulated algorithmic grid yields it analytically.

\subsubsection{\texorpdfstring{6.3 Deriving the Scattering Amplitude
(The \(S\)-Matrix
Ghost)}{6.3 Deriving the Scattering Amplitude (The S-Matrix Ghost)}}\label{deriving-the-scattering-amplitude-the-s-matrix-ghost}

Mainstream physics asserts particle interactions require continuous
fields, virtual particle propagators, and continuum Riemann integrals
(Feynman diagram loops over the \(S\)-matrix).

In \(\ell^1\), a particle is solely a localized Regge defect (a missing
or surplus triangle disrupting the \(\{3,6\}\) geometric equilibrium).
When two defects converge geodedically on the graph, their topological
boundaries overlap, locally spiking the \(\ell^1\) sum of
inconsistencies. The healing algorithm is forced to resolve the combined
defect coordinate set to iteratively minimize the global \(\ell^1\)
action. However, because the graph is strictly discrete, there are only
a strictly finite, calculable number of geometric tiling pathways to
fracture and re-tile the boundary into outgoing defect configurations.
The continuum Path Integral is simply the \(\ell^2\) continuous
hallucination of a discrete graph routing optimization problem. The
physical ``Scattering Amplitude'' is precisely the un-normalized sum of
the finite algebraic geometric permutations necessary for the 2-simplex
to heal a local tear.

\begin{center}\rule{0.5\linewidth}{0.5pt}\end{center}

\subsection{7. Dark Matter from Non-Resonant
Defects}\label{dark-matter-from-non-resonant-defects}

\subsubsection{7.1 Resonant and Non-Resonant
Configurations}\label{resonant-and-non-resonant-configurations}

In the \(\ell^1\) framework, Standard Model particles are \(\ell^1\)
defect configurations that exactly match the gauge representation
requirements: quarks carry \((3, 2, Y)\) under
\(SU(3) \times SU(2) \times U(1)\); leptons carry \((1, 2, Y)\) or
\((1, 1, Y)\). These are \emph{resonant} defects: they fit the simplex
boundary structure.

Not every \(\ell^1\) defect configuration factorizes into Standard Model
representations. Configurations that fail to match any \((R_3, R_2, Y)\)
assignment cannot couple to the Standard Model gauge fields. These
\emph{non-resonant} defects have four properties:

\begin{enumerate}
\def\labelenumi{\arabic{enumi}.}
\tightlist
\item
  \textbf{Massive}: they carry \(\ell^1\) defect energy.
\item
  \textbf{Invisible}: they cannot emit or absorb Standard Model gauge
  bosons.
\item
  \textbf{Gravitating}: they contribute to the Regge deficit.
\item
  \textbf{Stable}: no Standard Model decay channel is available.
\end{enumerate}

These properties are qualitatively consistent with those attributed to
dark matter.

\subsubsection{7.2 Topological Estimate}\label{topological-estimate}

On the single 2-simplex (3 edges, each carrying \(\pm 1\) orientation),
there are \(2^3 = 8\) topological configurations. Of these, 4 have
holonomy \(+1\) (consistent loop, resonant) and 4 have holonomy \(-1\)
(frustrated loop, non-resonant). The non-resonant fraction is
\(4/8 = 0.50\) (\texttt{execute\_50}).

The observed dark matter fraction is
\(\Omega_{\mathrm{DM}} / (\Omega_{\mathrm{DM}} + \Omega_b) = 0.27 / (0.27 + 0.05) = 0.84\).

The combinatorial estimate demonstrates that non-resonant configurations
are generically abundant. The precise cosmological abundance depends on
global lattice dynamics and is not predicted by the single-simplex
model.

\begin{center}\rule{0.5\linewidth}{0.5pt}\end{center}

\subsection{8. Dark Energy from Subdivision
Overshoot}\label{dark-energy-from-subdivision-overshoot}

\subsubsection{8.1 The Mechanism}\label{the-mechanism}

The \(\ell^1\) healing algorithm inserts vertices to reduce tension. The
equilibrium is the \(\{3,6\}\) tiling with valence 6. But the algorithm
is local: it cannot globally coordinate to produce exactly valence 6
everywhere.

In regions where subdivision overshoots, the local valence exceeds 6:

\[\text{valence} = 6 + \varepsilon, \quad \varepsilon > 0\]

The Regge deficit becomes negative:
\(\epsilon = 2\pi - (6 + \varepsilon) \times \pi/3 = -\varepsilon\pi/3\).
Negative curvature produces a repulsive gravitational effect, driving
accelerated expansion.

\subsubsection{8.2 The Cosmological Constant
Problem}\label{the-cosmological-constant-problem}

In quantum field theory, the cosmological constant problem asks: why is
\(\Lambda \sim 10^{-122} M_P^4\) instead of \(\sim M_P^4\)? In the
\(\ell^1\) framework:

\begin{itemize}
\tightlist
\item
  \(\Lambda = 0\) would require perfect global coordination of the
  subdivision algorithm, which is impossible for a local process.
\item
  \(\Lambda \sim M_P^4\) would require every vertex to have valence 7, a
  catastrophic overshoot inconsistent with the near-flatness of the
  observed universe.
\item
  \(\Lambda \sim 10^{-122} M_P^4\) means the local algorithm is almost
  perfect: approximately 1 extra triangle per \(10^{61}\) vertices.
\end{itemize}

The smallness of \(\Lambda\) is not fine-tuning. It reflects the
efficiency of the \(\ell^1\) healing algorithm: a local process that is
\(99.99\ldots\%\) correct globally (\texttt{execute\_51}).

These interpretations are qualitative and are not presented as derived
predictions. A quantitative prediction of \(\Lambda\) and dark matter
mass spectra from \(\ell^1\) dynamics remains an open problem.

\begin{center}\rule{0.5\linewidth}{0.5pt}\end{center}

\subsection{\texorpdfstring{9. The Norm Tower:
\(\ell^1 \to \ell^2 \to \ell^\infty\)}{9. The Norm Tower: \textbackslash ell\^{}1 \textbackslash to \textbackslash ell\^{}2 \textbackslash to \textbackslash ell\^{}\textbackslash infty}}\label{the-norm-tower-ell1-to-ell2-to-ellinfty}

\subsubsection{9.1 Three Norms, Three
Scales}\label{three-norms-three-scales}

Each \(\ell^p\) norm describes physics at a characteristic scale:

\begin{longtable}[]{@{}
  >{\raggedright\arraybackslash}p{(\linewidth - 6\tabcolsep) * \real{0.1579}}
  >{\raggedright\arraybackslash}p{(\linewidth - 6\tabcolsep) * \real{0.1842}}
  >{\raggedright\arraybackslash}p{(\linewidth - 6\tabcolsep) * \real{0.2368}}
  >{\raggedright\arraybackslash}p{(\linewidth - 6\tabcolsep) * \real{0.4211}}@{}}
\toprule\noalign{}
\begin{minipage}[b]{\linewidth}\raggedright
Norm
\end{minipage} & \begin{minipage}[b]{\linewidth}\raggedright
Scale
\end{minipage} & \begin{minipage}[b]{\linewidth}\raggedright
Physics
\end{minipage} & \begin{minipage}[b]{\linewidth}\raggedright
Characteristic
\end{minipage} \\
\midrule\noalign{}
\endhead
\bottomrule\noalign{}
\endlastfoot
\(\ell^1\) & Planck & Discrete defect counting & No cancellations;
monotone correction \\
\(\ell^2\) & Quantum & Hilbert space, Born rule & Interference allowed;
Parseval's identity \\
\(\ell^\infty\) & Classical & Deterministic trajectories & Maximum
defect dominates \\
\end{longtable}

\subsubsection{9.2 The Norm Inequality}\label{the-norm-inequality}

For any vector \(x\): \(\|x\|_\infty \leq \|x\|_2 \leq \|x\|_1\). This
is verified for all test configurations at \texttt{execute\_52}. The
inequality is strict when more than one component is nonzero. The
\(\ell^1\) norm is maximally sensitive (it counts all defects
individually); the \(\ell^\infty\) norm sees only the largest; the
\(\ell^2\) norm interpolates.

\subsubsection{9.3 The Emergence Chain}\label{the-emergence-chain}

The scaling hierarchy is:

\[\ell^1 \text{ (Planck)} \xrightarrow{\text{Parseval}} \ell^2 \text{ (quantum)} \xrightarrow{\text{law of large numbers}} \ell^\infty \text{ (classical)} \xrightarrow{\text{Regge limit}} \text{GR (cosmological)}\]

Each level is a coarse-graining of the one above it. Information is lost
at each step:

\begin{enumerate}
\def\labelenumi{\arabic{enumi}.}
\item
  \textbf{\(\ell^1 \to \ell^2\)}: The discrete Fourier transform and
  Parseval's identity (Paper 008, Section 4) transform the \(\ell^1\)
  defect counting into \(\ell^2\) Hilbert space structure. This aligns
  with the phase-averaging derivation of the Born rule (Paper 007a).
\item
  \textbf{\(\ell^2 \to \ell^\infty\)}: The law of large numbers applied
  to many quantum measurements. In the limit of many particles, the
  dominant eigenvalue controls the dynamics. This is the classical
  limit.
\item
  \textbf{\(\ell^\infty \to\) GR}: The Regge limit of the classical
  tiling. Smooth curvature emerges from the statistical aggregate of
  discrete deficits. This is general relativity.
\end{enumerate}

The \(\ell^1\) lattice is the only exact description. Every other level
of physics is an approximation obtained by changing the norm.

\begin{center}\rule{0.5\linewidth}{0.5pt}\end{center}

\subsection{10. The Derivation
Structure}\label{the-derivation-structure}

\subsubsection{10.1 Overview}\label{overview}

The framework proposes a constrained derivation chain beginning from:

\begin{itemize}
\tightlist
\item
  a minimal inconsistency (sign distinction)
\item
  a monotone correction principle
\end{itemize}

and generating a hierarchy of structures through successive constraints.

Each step is either:

\begin{itemize}
\tightlist
\item
  \textbf{derived} (structurally forced)
\item
  \textbf{constructed} (given prior results)
\item
  \textbf{interpreted} (mapping to known physics)
\end{itemize}

\subsubsection{10.2 Structural
Progression}\label{structural-progression}

The derivation proceeds in stages:

\paragraph{Stage I: Norm and Discrete
Structure}\label{stage-i-norm-and-discrete-structure}

\begin{enumerate}
\def\labelenumi{\arabic{enumi}.}
\item
  \textbf{\(\ell^1\) norm selection} Forced by additivity, locality, and
  monotone correction
\item
  \textbf{Emergence of \(N = 3\) structure} From minimal resolution of
  sign inconsistency
\item
  \textbf{Discrete Fourier structure (DFT)} Representation theory of
  \(\mathbb{Z}_3\)
\item
  \textbf{Hilbert structure and Born rule} Via Parseval correspondence
  under transform
\end{enumerate}

\paragraph{Stage II: Operator and Measure
Constraints}\label{stage-ii-operator-and-measure-constraints}

\begin{enumerate}
\def\labelenumi{\arabic{enumi}.}
\setcounter{enumi}{4}
\item
  \textbf{Transport amplitude structure (\(a = 3/\pi\))} From symmetry,
  normalization, and domain constraints
\item
  \textbf{Integral representation (RMK)} Operator expressed via measure,
  observable, domain
\item
  \textbf{Unitarity deviation parameter} Derived from structural
  constraints on operators
\end{enumerate}

\paragraph{Stage III: Gauge Structure}\label{stage-iii-gauge-structure}

\begin{enumerate}
\def\labelenumi{\arabic{enumi}.}
\setcounter{enumi}{7}
\tightlist
\item
  \textbf{\(SU(3)\) from closed loop structure}
\item
  \textbf{\(SU(2)\) from bifurcation symmetry}
\item
  \textbf{\(U(1)\) from edge-level transport}
\end{enumerate}

These arise as:

\begin{quote}
distinct dimensional projections of the same underlying structure
\end{quote}

\paragraph{Stage IV: Phenomenological
Correspondence}\label{stage-iv-phenomenological-correspondence}

\begin{enumerate}
\def\labelenumi{\arabic{enumi}.}
\setcounter{enumi}{10}
\tightlist
\item
  \textbf{Coupling relationships (GUT scale)}
\item
  \textbf{Generation count (\(N_g = 3\))}
\item
  \textbf{Electroweak structure and mixing angle}
\end{enumerate}

These steps involve:

\begin{itemize}
\tightlist
\item
  structural derivation
\item
  plus phenomenological anchoring
\end{itemize}

\paragraph{Stage V: Geometric
Extension}\label{stage-v-geometric-extension}

\begin{enumerate}
\def\labelenumi{\arabic{enumi}.}
\setcounter{enumi}{13}
\tightlist
\item
  \textbf{Discrete tiling (\(\{3,6\}\))}
\item
  \textbf{Curvature from Regge defects (Section 3)}
\item
  \textbf{Planck scale as defect threshold}
\end{enumerate}

\paragraph{Stage VI: Structural
Consequences}\label{stage-vi-structural-consequences}

\begin{enumerate}
\def\labelenumi{\arabic{enumi}.}
\setcounter{enumi}{16}
\tightlist
\item
  \textbf{Non-invertibility of coarse-graining (Section 5)}
\item
  \textbf{Norm hierarchy \(\ell^1 \to \ell^2 \to \ell^\infty\)}
\item
  \textbf{Qualitative dark sector mechanisms}
\end{enumerate}

\subsubsection{10.3 Nature of the
Derivation}\label{nature-of-the-derivation}

The chain is not purely deductive in the classical axiomatic sense.

Instead, it is:

\begin{quote}
a constrained construction in which each stage reduces available degrees
of freedom
\end{quote}

Key characteristics:

\begin{itemize}
\tightlist
\item
  No continuous free parameters introduced prior to phenomenological
  matching
\item
  Structural constraints propagate forward
\item
  Later stages depend only on earlier derived structure
\end{itemize}

\subsubsection{10.4 Boundary of Validity}\label{boundary-of-validity}

The framework distinguishes clearly between:

\paragraph{Structurally Derived}\label{structurally-derived}

\begin{itemize}
\tightlist
\item
  discrete lattice structure
\item
  norm hierarchy
\item
  gauge decomposition
\item
  curvature as defect
\item
  Newton's constant \(G\)
\item
  Einstein Field Equations
\end{itemize}

\paragraph{Partially Derived /
Anchored}\label{partially-derived-anchored}

\begin{itemize}
\tightlist
\item
  coupling constants
\item
  energy scales
\item
  mixing parameters
\end{itemize}

\paragraph{Open}\label{open}

\begin{itemize}
\tightlist
\item
  cosmological constant (quantitative)
\item
  fermion mass spectrum
\end{itemize}

\subsubsection{10.5 Interpretation}\label{interpretation-1}

The derivation suggests that:

\begin{quote}
a large portion of known physical structure may arise from constraints
on resolving minimal inconsistency
\end{quote}

However:

\begin{itemize}
\tightlist
\item
  the framework is not yet complete
\item
  several key quantitative components remain external
\end{itemize}

\subsubsection{10.6 Computational Role}\label{computational-role}

The associated computational pipeline:

\begin{itemize}
\tightlist
\item
  verifies internal consistency
\item
  confirms constraint propagation
\item
  checks absence of contradictions
\end{itemize}

It does \textbf{not} constitute experimental validation.

\begin{center}\rule{0.5\linewidth}{0.5pt}\end{center}

\subsection{\texorpdfstring{11. Operator Uniqueness Under \(\ell^1\)
Structural
Constraints}{11. Operator Uniqueness Under \textbackslash ell\^{}1 Structural Constraints}}\label{operator-uniqueness-under-ell1-structural-constraints}

We formalize the uniqueness of the transport amplitude within the
admissible operator class defined by the \(\ell^1\) framework.

\subsubsection{11.1 Problem Statement}\label{problem-statement}

Let an admissible transport operator be defined as a functional:

\[
\mathcal{T}(f) = \int_{\Omega} f(\theta) \, d\mu(\theta)
\]

where:

\begin{itemize}
\tightlist
\item
  \(\mu\) is a measure on \(S^1\)
\item
  \(f\) is a real-valued observable
\item
  \(\Omega \subset S^1\) is a domain determined by local topology
\end{itemize}

We seek to characterize all triples \((\mu, f, \Omega)\) satisfying the
structural constraints imposed by the \(\ell^1\) causal lattice.

\subsubsection{11.2 Structural
Constraints}\label{structural-constraints}

The admissible operator must satisfy:

\paragraph{(C1) Shift Invariance}\label{c1-shift-invariance}

\[
\mu(\theta + \phi) = \mu(\theta)
\] for all \(\phi \in S^1\)

→ implies \(\mu\) is proportional to Haar measure

\paragraph{(C2) Linearity}\label{c2-linearity}

\[
\mathcal{T}(a f_1 + b f_2) = a \mathcal{T}(f_1) + b \mathcal{T}(f_2)
\]

→ restricts operator to integral form

\paragraph{(C3) Reality}\label{c3-reality}

\[
f(\theta) \in \mathbb{R}
\]

\paragraph{(C4) Normalization}\label{c4-normalization}

\[
f(0) = 1
\]

\paragraph{\texorpdfstring{(C5) Even Symmetry (\(\ell^1\) Edge
Invariance)}{(C5) Even Symmetry (\textbackslash ell\^{}1 Edge Invariance)}}\label{c5-even-symmetry-ell1-edge-invariance}

\[
f(\theta) = f(-\theta)
\]

→ removes odd components

\paragraph{(C6) Locality / Single-Step
Transport}\label{c6-locality-single-step-transport}

Operator depends only on a single-edge interaction

→ excludes higher harmonics \(f(n\theta)\), for \(n > 1\)

\paragraph{(C7) Domain Constraint (Discrete
Topology)}\label{c7-domain-constraint-discrete-topology}

The domain \(\Omega\) is determined by:

\begin{itemize}
\tightlist
\item
  cycle graph \(C_3\)
\item
  vertex degree \(d = 2\)
\item
  equal Voronoi partition
\end{itemize}

\[
\Omega = [-\pi/6, \pi/6]
\]

\subsubsection{11.3 Classification of Observable
Space}\label{classification-of-observable-space}

The general real-valued function on \(S^1\) can be written:

\[
f(\theta) = a \cos\theta + b \sin\theta
\]

Applying constraints:

\begin{itemize}
\tightlist
\item
  Even symmetry ⇒ \(b = 0\)
\item
  Normalization ⇒ \(a = 1\)
\end{itemize}

Thus:

\[
f(\theta) = \cos\theta
\]

This is the \textbf{unique admissible observable}.

\subsubsection{11.4 Measure Uniqueness}\label{measure-uniqueness}

By (C1), the measure is uniquely:

\[
d\mu = \frac{d\theta}{|\Omega|}
\]

(Haar measure normalized on the domain)

\subsubsection{11.5 Domain Uniqueness}\label{domain-uniqueness}

The domain \(\Omega\) is fixed by:

\begin{itemize}
\tightlist
\item
  minimal closed topology (\(N = 3\))
\item
  vertex degree (\(d = 2\))
\item
  equal partition of angular space
\end{itemize}

Thus:

\[
\Omega = [-\pi/6, \pi/6]
\]

No continuous degree of freedom remains.

\subsubsection{11.6 Result}\label{result}

The transport amplitude is therefore uniquely:

\[
\mathcal{T} = \frac{1}{|\Omega|} \int_{-\pi/6}^{\pi/6} \cos\theta \, d\theta
= \frac{\sin(\pi/6)}{\pi/6}
= \frac{3}{\pi}
\]

\subsubsection{11.7 Uniqueness Statement}\label{uniqueness-statement}

Under constraints (C1--C7):

\begin{quote}
Within the constraint class (C1--C7), there exists a unique admissible
transport operator. We do not claim that the constraint class itself is
uniquely forced.
\end{quote}

Equivalently:

\begin{itemize}
\tightlist
\item
  Measure space: dimension \(0\) (unique)
\item
  Observable space: dimension \(0\) (unique)
\item
  Domain: dimension \(0\) (unique)
\end{itemize}

Thus:

\[
\dim(\text{operator space}) = 0
\]

and the transport amplitude \(3/\pi\) is the \textbf{unique output} of
the construction.

\subsubsection{11.8 Exclusion of
Alternatives}\label{exclusion-of-alternatives}

All alternative constructions violate at least one constraint:

\begin{longtable}[]{@{}ll@{}}
\toprule\noalign{}
Alternative & Violation \\
\midrule\noalign{}
\endhead
\bottomrule\noalign{}
\endlastfoot
Non-Haar measure & breaks shift invariance \\
\(\sin\theta\) component & breaks even symmetry \\
\(\cos(n\theta), n>1\) & violates locality / single-step \\
Larger domain & violates discrete topology \\
Weighted domain & breaks symmetry \\
\end{longtable}

\subsubsection{11.9 Interpretation}\label{interpretation-2}

The value \(3/\pi\) is not introduced or fitted.

It arises as:

\begin{quote}
the unique invariant amplitude of single-step transport on a discrete
\(\ell^1\)-stabilized 2-simplex under symmetry, locality, and
normalization constraints.
\end{quote}

The derivation of the transport amplitude is conditional on the discrete
topology and locality constraints derived in earlier sections.
Alternative topologies or nonlocal operators may yield different
amplitudes, but such constructions fall outside the present framework.

\begin{center}\rule{0.5\linewidth}{0.5pt}\end{center}

\subsection{\texorpdfstring{12. Gauge Structure from \(\ell^1\) Boundary
Constraints}{12. Gauge Structure from \textbackslash ell\^{}1 Boundary Constraints}}\label{gauge-structure-from-ell1-boundary-constraints}

We formalize the emergence and uniqueness of the gauge structure within
the \(\ell^1\) causal lattice.

\subsubsection{12.1 Problem Statement}\label{problem-statement-1}

Given:

\begin{itemize}
\tightlist
\item
  a discrete \(\ell^1\)-stabilized 2-simplex,
\item
  with minimal closed topology (\(N = 3\)),
\item
  and transport defined along edges, hinges, and closed loops,
\end{itemize}

we classify all admissible symmetry structures compatible with:

\begin{itemize}
\tightlist
\item
  locality,
\item
  unitary transport,
\item
  closure under composition,
\item
  and \(\ell^1\) boundary constraints.
\end{itemize}

\subsubsection{12.2 Structural
Decomposition}\label{structural-decomposition}

The 2-simplex admits three distinct geometric structures:

\begin{longtable}[]{@{}lll@{}}
\toprule\noalign{}
Structure & Dimension & Interpretation \\
\midrule\noalign{}
\endhead
\bottomrule\noalign{}
\endlastfoot
Vertex set & 0D & states \\
Edge & 1D & transport \\
Open path (2 edges) & 2D boundary & bifurcation \\
Closed loop (3 edges) & 2D closed boundary & holonomy \\
\end{longtable}

These correspond to \textbf{three distinct operator domains}.

\subsubsection{12.3 Constraint Set}\label{constraint-set}

Any admissible symmetry group \(G\) must satisfy:

\paragraph{(G1) Locality}\label{g1-locality}

Transformations act on single-step transport (edges or minimal loops)

\paragraph{(G2) Unitarity}\label{g2-unitarity}

Transport preserves norm under composition:

\[
U^\dagger U = I
\]

\paragraph{(G3) Closure}\label{g3-closure}

Group structure must be closed under composition of local transport
operations

\paragraph{(G4) Minimality}\label{g4-minimality}

No higher-dimensional cell structure beyond the 2-simplex is permitted

→ excludes \(N > 3\) polygonal structures

\paragraph{(G5) Boundary Compatibility}\label{g5-boundary-compatibility}

Symmetry must respect:

\begin{itemize}
\tightlist
\item
  edge transport (1D)
\item
  hinge transitions (2D open)
\item
  loop holonomy (2D closed)
\end{itemize}

\subsubsection{12.4 Classification by
Dimension}\label{classification-by-dimension}

We now classify admissible symmetry structures by boundary type.

\paragraph{(A) Closed Loop Symmetry (3-edge
holonomy)}\label{a-closed-loop-symmetry-3-edge-holonomy}

The minimal closed structure is:

\[
C_3
\]

Transport around the loop defines:

\[
U = \exp(i \theta_a T^a)
\]

Constraints:

\begin{itemize}
\tightlist
\item
  3 independent directions
\item
  non-abelian closure required
\item
  unitary representation
\end{itemize}

The unique compact simple Lie group satisfying:

\begin{itemize}
\tightlist
\item
  dimension = 3 generators (fundamental representation dimension 3)
\item
  non-abelian closure
\item
  compatibility with complex vector space \(\mathbb{C}^3\)
\end{itemize}

is:

\[
SU(3)
\]

\paragraph{(B) Open Boundary Symmetry (2-edge
hinge)}\label{b-open-boundary-symmetry-2-edge-hinge}

The minimal open structure is a bifurcation:

\begin{itemize}
\tightlist
\item
  2 connected edges
\item
  no closure constraint
\end{itemize}

Transport must:

\begin{itemize}
\tightlist
\item
  be unitary
\item
  allow mixing between two states
\item
  preserve normalization
\end{itemize}

This uniquely selects:

\[
SU(2)
\]

as the group of unitary transformations on \(\mathbb{C}^2\)

\paragraph{(C) Edge Transport (1D)}\label{c-edge-transport-1d}

A single edge supports:

\begin{itemize}
\tightlist
\item
  phase transport only
\item
  no internal mixing
\end{itemize}

Thus:

\[
U = e^{i\theta}
\]

This defines:

\[
U(1)
\]

\subsubsection{12.5 Exclusion of Higher
Groups}\label{exclusion-of-higher-groups}

We now show that larger groups are inadmissible.

\paragraph{\texorpdfstring{(E1)
\(SU(N), N > 3\)}{(E1) SU(N), N \textgreater{} 3}}\label{e1-sun-n-3}

Requires:

\begin{itemize}
\tightlist
\item
  \(N\)-gon minimal closed structure
\item
  \(N > 3\) edges per loop
\end{itemize}

Violates:

\begin{itemize}
\tightlist
\item
  (G4) minimality
\item
  \(\ell^1\) transport optimality (increased path length)
\end{itemize}

\paragraph{\texorpdfstring{(E2) \(SO(N)\)}{(E2) SO(N)}}\label{e2-son}

Acts on real vector spaces.

Violates:

\begin{itemize}
\tightlist
\item
  complex structure required by DFT / Hilbert emergence
\end{itemize}

\paragraph{(E3) Product Extensions}\label{e3-product-extensions}

Additional factors require:

\begin{itemize}
\tightlist
\item
  independent degrees of freedom
\item
  additional geometric structure
\end{itemize}

Not present in 2-simplex topology

\paragraph{(E4) Exceptional Groups}\label{e4-exceptional-groups}

Require:

\begin{itemize}
\tightlist
\item
  higher-dimensional algebraic structure
\item
  no embedding in 2-simplex boundary
\end{itemize}

\subsubsection{12.6 Result}\label{result-1}

The admissible symmetry decomposition is uniquely:

\[
SU(3) \times SU(2) \times U(1)
\]

with:

\begin{itemize}
\tightlist
\item
  \(SU(3)\): closed loop holonomy
\item
  \(SU(2)\): open boundary mixing
\item
  \(U(1)\): edge phase transport
\end{itemize}

The identification of \(SU(3)\), \(SU(2)\), and \(U(1)\) should be
understood as the minimal compact unitary groups compatible with the
dimensionality and compositional structure of the boundary operators. We
do not exclude the existence of alternative representations without
additional constraints on representation faithfulness and
irreducibility.

\subsubsection{12.7 Uniqueness Statement}\label{uniqueness-statement-1}

Under constraints (G1--G5):

\begin{quote}
Within the constraint class (G1--G5), there exists a unique admissible
gauge structure compatible with the \(\ell^1\)-stabilized 2-simplex. We
do not claim that the constraint class itself is uniquely forced.
\end{quote}

Equivalently:

\begin{itemize}
\tightlist
\item
  Closed symmetry space: dimension 0 (unique)
\item
  Open symmetry space: dimension 0 (unique)
\item
  Edge symmetry space: dimension 0 (unique)
\end{itemize}

Thus:

\[
\dim(\text{gauge structure space}) = 0
\]

\subsubsection{12.8 Interpretation}\label{interpretation-3}

The Standard Model gauge group is not:

\begin{itemize}
\tightlist
\item
  postulated,
\item
  or selected from a larger symmetry,
\end{itemize}

but arises as:

\begin{quote}
the unique decomposition of transport symmetries across the dimensional
boundaries of the minimal \(\ell^1\)-stabilized topology.
\end{quote}

\begin{center}\rule{0.5\linewidth}{0.5pt}\end{center}

\subsection{13. Zero Free Parameter
Theorem}\label{zero-free-parameter-theorem}

We now combine the preceding results to establish the structural
rigidity of the \(\ell^1\) causal lattice framework.

\subsubsection{13.1 Statement of the
Problem}\label{statement-of-the-problem}

A physical theory is said to possess \textbf{free parameters} if:

\begin{itemize}
\tightlist
\item
  continuous degrees of freedom must be externally specified, or
\item
  multiple inequivalent constructions satisfy the same foundational
  constraints.
\end{itemize}

We examine whether the \(\ell^1\) causal lattice admits any such
freedom.

\subsubsection{13.2 Structural Components}\label{structural-components}

The construction proceeds through the following elements:

\begin{longtable}[]{@{}lll@{}}
\toprule\noalign{}
Component & Role & Status \\
\midrule\noalign{}
\endhead
\bottomrule\noalign{}
\endlastfoot
\(\ell^1\) norm & defect measurement & forced \\
2-simplex (\(N = 3\)) & minimal resolution topology & forced \\
Fourier structure & spectral representation & forced \\
Operator (transport amplitude) & interaction strength & unique \\
Gauge structure & symmetry decomposition & unique \\
Domain geometry & angular constraint & unique \\
\end{longtable}

Each component is determined by prior constraints.

\subsubsection{13.3 Degree-of-Freedom
Accounting}\label{degree-of-freedom-accounting}

We count remaining degrees of freedom:

\paragraph{(A) Measure space}\label{a-measure-space}

By Haar invariance:

\[
\dim(\mu) = 0
\]

\paragraph{(B) Observable space}\label{b-observable-space}

After symmetry and normalization:

\[
\dim(f) = 0
\]

\paragraph{(C) Domain}\label{c-domain}

From discrete topology:

\[
\dim(\Omega) = 0
\]

\paragraph{(D) Operator}\label{d-operator}

Constructed from (A--C):

\[
\dim(\mathcal{T}) = 0
\]

\paragraph{(E) Gauge structure}\label{e-gauge-structure}

From boundary classification:

\[
\dim(G) = 0
\]

\subsubsection{13.4 Result}\label{result-2}

The total number of internal free parameters is:

\[
\dim(\text{theory}) = 0
\]

\subsubsection{13.5 Interpretation}\label{interpretation-4}

All primary structures arise from:

\begin{itemize}
\tightlist
\item
  discrete topology
\item
  symmetry constraints
\item
  locality and minimality
\end{itemize}

No continuous parameter is introduced prior to phenomenological
matching.

Thus:

\begin{quote}
The \(\ell^1\) causal lattice defines a \textbf{framework with no
continuous free parameters within the specified constraint class} at the
level of its internal construction.
\end{quote}

\subsubsection{13.6 Relation to Physical
Constants}\label{relation-to-physical-constants}

Observed physical constants enter in two ways:

\paragraph{(1) Structurally determined}\label{structurally-determined}

Examples:

\begin{itemize}
\tightlist
\item
  transport amplitude (\(3/\pi\))
\item
  gauge decomposition
\item
  generation count
\end{itemize}

These arise from zero-parameter construction.

\paragraph{(2) Phenomenologically
anchored}\label{phenomenologically-anchored}

Examples:

\begin{itemize}
\tightlist
\item
  cosmological constant (\(\Lambda\))
\item
  fermion mass spectrum
\end{itemize}

These correspond to:

\begin{quote}
quantities not yet derived from the discrete structure
\end{quote}

and are explicitly identified as open problems.

\subsubsection{13.7 Consequence}\label{consequence}

The framework implies:

\begin{quote}
any successful geometric or symmetry extension must preserve zero
internal degrees of freedom
\end{quote}

or identify precisely where additional topological structure enters.

\subsubsection{13.8 Falsifiability}\label{falsifiability}

The zero-parameter nature introduces a strong constraint:

\begin{quote}
any alternative construction satisfying the same structural axioms must
reproduce the same outputs.
\end{quote}

Thus, the framework is falsified if:

\begin{itemize}
\tightlist
\item
  an admissible operator exists with amplitude \(\neq 3/\pi\), or
\item
  an admissible symmetry structure differs from
  \(SU(3) \times SU(2) \times U(1)\), or
\item
  a consistent \(\ell^1\) construction produces additional free
  parameters
\end{itemize}

\subsubsection{13.9 Summary}\label{summary}

The \(\ell^1\) causal lattice is not:

\begin{itemize}
\tightlist
\item
  a parameterized model,
\item
  nor a family of theories
\end{itemize}

but:

\begin{quote}
a single constrained construction with no internal degrees of freedom.
The only remaining freedom lies in the existence of the initial
inconsistency itself, which is not a parameter of the theory but the
condition under which any theory becomes possible.
\end{quote}

\subsubsection{13.10 On the Scope of Uniqueness
Claims}\label{on-the-scope-of-uniqueness-claims}

The uniqueness results established in Sections 11--13 are conditional on
the specified constraint set:

\begin{itemize}
\tightlist
\item
  locality (single-step transport)
\item
  linearity
\item
  symmetry (shift invariance, parity)
\item
  minimal topology (2-simplex)
\end{itemize}

A potential objection is that alternative admissible constructions may
exist outside this constraint set.

This objection is valid in principle. However:

\begin{quote}
any such construction must explicitly identify which constraint is
relaxed and demonstrate that the resulting structure preserves the
intended physical invariants.
\end{quote}

In particular, alternatives must account for:

\begin{itemize}
\tightlist
\item
  stability under composition
\item
  bounded transport
\item
  compatibility with discrete topology
\end{itemize}

Absent such a construction, the classification remains complete within
the stated domain.

\begin{center}\rule{0.5\linewidth}{0.5pt}\end{center}

\subsection{14. Open Problems}\label{open-problems}

\begin{enumerate}
\def\labelenumi{\arabic{enumi}.}
\item
  \textbf{Derive the cosmological constant quantitatively.} The
  subdivision overshoot mechanism predicts \(\Lambda > 0\) and explains
  its smallness qualitatively. A quantitative prediction requires a
  theory of the average overshoot rate.
\item
  \textbf{Derive the dark matter mass spectrum.} The non-resonant defect
  mechanism predicts massive, invisible, gravitating, stable particles.
  The mass distribution requires understanding which non-resonant
  configurations are dynamically accessible.
\item
  \textbf{Formalize the coarse-graining steps.} The transitions
  \(\ell^1 \to \ell^2\) (Parseval) and \(\ell^2 \to \ell^\infty\) (law
  of large numbers) are identified but not rigorously proven to be the
  unique coarse-graining maps preserving the relevant physical
  structure.
\item
  \textbf{Derive fermion masses and mixing.} The three generations are
  forced (Paper 009), but the mass hierarchy within and between
  generations remains underived.
\end{enumerate}

\begin{center}\rule{0.5\linewidth}{0.5pt}\end{center}

\subsection{15. Conclusion}\label{conclusion}

This work proposes that a substantial portion of known physical
structure can be derived from a single organizing principle: monotone
correction of local inconsistency under the \(\ell^1\) coboundary norm.

Starting from this principle and a minimal discrete seed, the framework
constructs a sequence of constrained structures, including:

\begin{itemize}
\tightlist
\item
  a discrete \(N = 3\) lattice and associated Fourier structure,
\item
  Hilbert space behavior via \(\ell^1 \to \ell^2\) correspondence,
\item
  a constrained transport amplitude and associated unitarity deviation,
\item
  a framework consistent with key structural features of the Standard
  Model gauge group \(SU(3) \times SU(2) \times U(1)\),
\item
  consistency with three fermion generations and electroweak structure,
\item
  discrete geometric tiling and curvature via Regge defects,
\item
  the Planck scale as a defect threshold,
\item
  and a hierarchy of effective descriptions arising from norm
  transitions.
\end{itemize}

Within this framework, several longstanding problems are reinterpreted:

\begin{itemize}
\tightlist
\item
  The hierarchy between gravitational and electroweak scales appears as
  a statement about lattice extent rather than parameter fine-tuning.
\item
  The cosmological constant reflects residual imbalance in a locally
  driven subdivision process rather than vacuum energy divergence.
\item
  Dark matter corresponds to defect configurations that do not couple to
  the Standard Model gauge structure.
\end{itemize}

The framework does not yet provide quantitative derivations of the
cosmological constant, a global dark matter mass spectrum, or fermion
mass hierarchies. These remain open problems.

The central claim is therefore not that all physical quantities have
been derived, but that:

\begin{quote}
a large fraction of known structure can be organized as the consequence
of a single constrained discrete mechanism, with remaining parameters
isolated and clearly identified.
\end{quote}

If correct, this shifts the role of unification:

\begin{quote}
not as the merger of independent theories, but as the identification of
a common discrete substrate from which multiple effective descriptions
emerge.
\end{quote}

\begin{center}\rule{0.5\linewidth}{0.5pt}\end{center}

\subsection{References}\label{references}

{[}1{]} T. Regge, ``General relativity without coordinates,'' \emph{Il
Nuovo Cimento} \textbf{19}, 558 (1961).

{[}2{]} J. Maldacena, ``The large-\(N\) limit of superconformal field
theories and supergravity,'' \emph{Int. J. Theor. Phys.} \textbf{38},
1113 (1999).

{[}3{]} C. Rovelli, \emph{Quantum Gravity} (Cambridge University Press,
2004).

{[}4{]} J. Polchinski, \emph{String Theory}, Vols. 1--2 (Cambridge
University Press, 1998).

{[}5{]} R. D. Sorkin, ``Causal sets: Discrete gravity,'' in
\emph{Lectures on Quantum Gravity}, eds.~A. Gomberoff and D. Marolf
(Springer, 2005).

{[}6{]} Particle Data Group (R. L. Workman et al.), ``Review of Particle
Physics,'' \emph{PTEP} \textbf{2022}, 083C01 (2022). Updated 2024.

{[}7{]} Planck Collaboration (N. Aghanim et al.), ``Planck 2018 results.
VI. Cosmological parameters,'' \emph{Astron. Astrophys.} \textbf{641},
A6 (2020).

{[}8{]} J. H. Carroll, ``Paper 000: Projection Obstruction Theory -
Retraction Nonlinearity, l1 Rigidity, and Density Scaling'' (2026).

{[}9{]} J. H. Carroll, ``Paper 001: The Free L1 Seminorm on Banach
Presheaf Coboundaries'' (2026).

{[}10{]} J. H. Carroll, ``Paper 002: Coordinate-Wise Additivity and the
L1 Norm on Finite Graph Cochains'' (2026).

{[}10{]} J. H. Carroll, ``Paper 003: Hodge Structure and Gauge
Equivalence of L1 Defect Fields'' (2026).

{[}11{]} J. H. Carroll, ``Paper 004: An L1 Obstruction Framework for
Coboundary Defects'' (2026).

{[}12{]} J. H. Carroll, ``Paper 005: Autopoietic Cohomology - Iterative
Obstruction Repair on Causal Posets'' (2026).

{[}13{]} J. H. Carroll, ``Paper 006: A Common L1 Defect Framework for
Quantum, Causal, and Gravitational Systems'' (2026).

{[}14{]} J. H. Carroll, ``Paper 007: Topological Constraint on the
Electromagnetic Coupling Scale from a Discrete \(\ell^1\) Obstruction''
(2026).

{[}15{]} J. H. Carroll, ``Paper 008: Emergent Gauge Symmetry from L1
Defect Flow'' (2026).

\begin{center}\rule{0.5\linewidth}{0.5pt}\end{center}

\end{document}
